%======================================================================
\chapter{Minimization algorithms}\label{app:opt}
%======================================================================
% Written for Chapter 7 (identification): the algorithms named there
% (Nelder-Mead, BFGS, Gauss-Newton, Levenberg-Marquardt) with their
% principles, costs and pitfalls. Code in python/examples/appF.

\framepart{Outline}
Chapter~\ref{ch:id} minimizes a cost $\mathcal J(\vp)$ over a few parameters,
each evaluation of $\mathcal J$ being a whole elastoplastic computation. This
appendix collects the algorithms it uses, in $\R^n$ with a smooth cost unless said
otherwise. They all repeat two decisions: a \emph{direction} in which the cost
decreases, and a \emph{step} along it, or a region in which a model of the cost is
trusted. They differ by the information they use: values only, the gradient, an
approximation of the Hessian built from gradients, or, for least squares, the
Jacobian of the residuals.

Section~\ref{sec:opt-descent} treats descent directions and line searches,
Section~\ref{sec:opt-newton} Newton and quasi-Newton methods,
Section~\ref{sec:opt-ls} least squares (Gauss--Newton, Levenberg--Marquardt),
Section~\ref{sec:opt-free} methods without derivatives, and
Section~\ref{sec:opt-bounds} bounds and scaling. The standard modern reference is
Nocedal and Wright~\cite{NocedalWright2006b}; Dennis and
Schnabel~\cite{DennisSchnabel1983b} and Fletcher~\cite{Fletcher1987} are
classics, Kelley~\cite{Kelley1999} is short and practical, and \emph{Numerical
Recipes}~\cite[Ch.~10, 15]{PressEtAl2007} gives readable codes and sound advice
for every method of this appendix. The equations $F(x)=0$ solved by Newton's
method are the subject of Appendix~\ref{app:iter}; here the unknown is a minimizer.

%======================================================================
\section{Descent directions and line searches}\label{sec:opt-descent}
\index{descent direction}\index{line search}%
At a point $\vp_k$ with the gradient $\vect g_k=\nabla\mathcal J(\vp_k)\neq\vzero$, a
direction $\vd_k$ is a \emph{descent direction} if $\vect g_k\cdot\vd_k<0$: a small
step along $\vd_k$ decreases $\mathcal J$. The iteration is
$\vp_{k+1}=\vp_k+a_k\vd_k$ with a step $a_k>0$. An exact minimization along the
line is wasteful; a \emph{line search} accepts any step that decreases $\mathcal J$
enough.
\begin{gbox}{Armijo and Wolfe conditions}\label{def:opt-wolfe}
\index{Armijo condition}\index{Wolfe conditions}%
With constants $0<c_1<c_2<1$ (typically $c_1=10^{-4}$, $c_2=0.9$), a step $a$ is
accepted if
\begin{equation}\label{eq:opt-wolfe}
\begin{aligned}
  &\mathcal J(\vp_k+a\vd_k)\le\mathcal J(\vp_k)+c_1\,a\,\vect g_k\cdot\vd_k
  &&\text{(sufficient decrease)},\\
  &\nabla\mathcal J(\vp_k+a\vd_k)\cdot\vd_k\ge c_2\,\vect g_k\cdot\vd_k
  &&\text{(curvature)}.
\end{aligned}
\end{equation}
Backtracking (try $a=1,\tfrac12,\tfrac14,\dots$) enforces the first
condition~\cite{Armijo1966}; the second excludes too short
steps~\cite{Wolfe1969}.
\end{gbox}
\noindent With such steps, any descent direction that does not become orthogonal
to the gradient gives $\nabla\mathcal J(\vp_k)\to\vzero$, for a cost bounded below
with a Lipschitz gradient~\cite[Sec.~3.2]{NocedalWright2006b}: the iteration
finds a stationary point. It may be a local minimum only, and nothing tells which.

\paragraph{Steepest descent is slow.}
\index{steepest descent}%
The obvious direction $\vd_k=-\vect g_k$ is the worst one in a valley. On the
quadratic $\tfrac12\vp\T\mat{A}\vp$ with exact line searches, the error decreases at
each iteration by the factor $(\kappa-1)/(\kappa+1)$, where $\kappa$ is the ratio of
the largest to the smallest eigenvalue of $\mat{A}$~\cite[Sec.~3.3]{NocedalWright2006b}.
The iterates zigzag across the valley. The Gauss--Newton matrix of the Norton
filament of Chapter~\ref{ch:id} has $\kappa=6\times10^3$: steepest descent would gain
about $3\times10^{-4}$ in the error per iteration. Every useful method corrects
the direction with curvature information.

%======================================================================
\section{Newton and quasi-Newton methods}\label{sec:opt-newton}
\index{Newton's method!minimization}\index{quasi-Newton method}%
Newton's method minimizes the quadratic model
$\mathcal J(\vp_k)+\vect g_k\cdot\vd+\tfrac12\vd\T\nabla^2\mathcal J_k\vd$:
$\nabla^2\mathcal J_k\,\vd_k=-\vect g_k$. It is Newton's method of
Appendix~\ref{app:iter} for the equation $\nabla\mathcal J=\vzero$, with quadratic
convergence near a minimum where the Hessian is positive definite. Far from it the
Hessian may be indefinite and $\vd_k$ not a descent direction; a line search, or a
modification of the Hessian, or a trust region (Section~\ref{sec:opt-ls}) is then
needed. In identification the Hessian is rarely available: it contains the
second derivatives of the whole computation with respect to the parameters.

\paragraph{BFGS.}
\index{BFGS method}%
Quasi-Newton methods build an approximation $\mat{H}_k$ of the inverse Hessian
from the gradients already computed. Along the last step
$\vs_k=\vp_{k+1}-\vp_k$ the gradient changed by $\vy_k=\vect g_{k+1}-\vect g_k$; a
Hessian should map $\vs_k$ to $\vy_k$, and the new approximation is required to do
so (the \emph{secant condition} $\mat{H}_{k+1}\vy_k=\vs_k$) while changing as
little as possible. The update of Broyden, Fletcher, Goldfarb and Shanno
(BFGS) is
\begin{equation}\label{eq:opt-bfgs}
  \mat{H}_{k+1}=\bigl(\mat{I}-\rho_k\vs_k\vy_k\T\bigr)\mat{H}_k\bigl(\mat{I}-\rho_k\vy_k\vs_k\T\bigr)
  +\rho_k\vs_k\vs_k\T,\qquad \rho_k=\frac1{\vy_k\cdot\vs_k},
\end{equation}
with the direction $\vd_k=-\mat{H}_k\vect g_k$. If $\vy_k\cdot\vs_k>0$, which the
Wolfe conditions guarantee, $\mat{H}_{k+1}$ stays symmetric positive definite and
$\vd_{k+1}$ is a descent direction. Near the minimum the convergence is
superlinear. Each iteration costs one gradient, the adjoint of Chapter~\ref{ch:id},
and a few products of vectors; the limited-memory version L-BFGS keeps only the
last pairs $(\vs_k,\vy_k)$ and works with millions of
parameters~\cite[Ch.~6--7]{NocedalWright2006b}.

%======================================================================
\section{Least squares: Gauss--Newton and Levenberg--Marquardt}\label{sec:opt-ls}
\index{least squares!nonlinear}\index{Gauss--Newton method}%
For $\mathcal J=\tfrac12\norm{\mat{r}(\vp)}^2$, the gradient and the Hessian are
\[
  \nabla\mathcal J=\mat{S}\T\mat{r},\qquad
  \nabla^2\mathcal J=\mat{S}\T\mat{S}+\sum_ir_i\nabla^2r_i,\qquad \mat{S}=\frac{\partial\mat{r}}{\partial\vp}.
\]
The first part of the Hessian needs only first derivatives, and the second is small
when the residuals are small or the model nearly linear. Dropping it gives the
\emph{Gauss--Newton} step
\begin{equation}\label{eq:opt-gn}
  \mat{S}_k\T\mat{S}_k\,\vd_k=-\mat{S}_k\T\mat{r}_k ,
\end{equation}
the solution of the linear least-squares problem
$\min_{\vd}\norm{\mat{r}_k+\mat{S}_k\vd}^2$, best computed by a QR factorization
of $\mat{S}_k$ rather than by the normal equations. The direction is a descent
direction whenever $\mat{S}_k$ has full rank. Convergence is quadratic for a
zero residual at the solution, linear otherwise, with a rate that grows with the
neglected term.

\paragraph{Levenberg--Marquardt.}
\index{Levenberg--Marquardt method}\index{trust region}%
When $\mat{S}\T\mat{S}$ is nearly singular, a flat valley, the Gauss--Newton step is
huge and meaningless. Levenberg~\cite{Levenberg1944} and
Marquardt~\cite{Marquardt1963} damp it:
\begin{equation}\label{eq:opt-lm}
  \bigl(\mat{S}_k\T\mat{S}_k+\nu_k\mat{D}_k\bigr)\vd_k=-\mat{S}_k\T\mat{r}_k ,
\end{equation}
with $\mat{D}_k$ the identity or the diagonal of $\mat{S}_k\T\mat{S}_k$. Large
$\nu_k$ gives a short step along the steepest descent; $\nu_k\to0$ gives
Gauss--Newton.
\begin{gbox}{The damped step is a trust-region step}\label{prop:opt-lm}
The step~\eqref{eq:opt-lm} with $\mat{D}=\mat{I}$ minimizes the quadratic model
$\tfrac12\norm{\mat{r}_k+\mat{S}_k\vd}^2$ on the ball $\norm{\vd}\le\Delta$, for the
radius $\Delta=\norm{\vd_k}$; conversely, every radius corresponds to a
damping $\nu\ge0$. The damping is updated by comparing the actual decrease of
$\mathcal J$ with the decrease predicted by the model.
\end{gbox}
\paragraph{Why it works.} Minimize $\tfrac12\norm{\mat{r}+\mat{S}\vd}^2$ under the
constraint $\tfrac12(\norm{\vd}^2-\Delta^2)\le0$. The Kuhn--Tucker conditions of
Section~\ref{sec:var-lagrange} give
$(\mat{S}\T\mat{S}+\nu\mat{I})\vd=-\mat{S}\T\mat{r}$ with the multiplier $\nu\ge0$,
zero if the Gauss--Newton step is inside the ball. The norm of the solution
decreases monotonically from the Gauss--Newton step ($\nu=0$) to zero
($\nu\to\infty$), so each radius has its damping. The damping is also a Tikhonov
regularization of the linearized problem: it penalizes the components of the step
along the small eigenvalues of $\mat{S}\T\mat{S}$, the poorly identified
combinations of parameters, and leaves the others.

In practice $\nu$ is decreased (by a factor $3$ to $10$) after a successful step
and increased after a failed one; implementations such as MINPACK~\cite{More1978}
control the radius directly, with the ratio of the actual to the predicted
decrease. Each iteration needs the residuals and the whole matrix $\mat{S}$, the
direct differentiation of Chapter~\ref{ch:id}, and solves a system of the size of
the number of parameters. For a few parameters and many observations it is the
method of choice.

%======================================================================
\section{Without derivatives}\label{sec:opt-free}
\index{derivative-free optimization}\index{Nelder--Mead method}%
When the gradient is unavailable, or the cost is noisy (an iterative direct solver
stopped at a loose tolerance), methods that use only values remain.
\begin{itemize}
  \item \emph{Finite-difference gradients} with a quasi-Newton method: $n$ or $2n$
  evaluations per gradient, with the step of Chapter~\ref{ch:id},
  Section~\ref{ssec:id-fd}. This is what most black-box calibration tools do.
  \item \emph{The simplex of Nelder and Mead}~\cite{NelderMead1965b} (Box~7.1): no
  model, only comparisons of values; robust to noise, efficient in two or three
  dimensions, without convergence guarantee beyond~\cite{LagariasEtAl1998b}, and
  easily stopped by curved valleys.
  \item \emph{Pattern search and model-based methods} (COBYLA, BOBYQA) build
  linear or quadratic models from interpolation points and come with convergence
  theory~\cite{ConnScheinbergVicente2009}.
  \item \emph{Global methods} (genetic and evolutionary algorithms, simulated
  annealing, Bayesian optimization on a surrogate) explore the whole admissible set
  and can escape local minima, at the price of hundreds to thousands of
  evaluations. They are best used to find a starting point for a local method.
\end{itemize}

%======================================================================
\section{Bounds and scaling}\label{sec:opt-bounds}
\index{scaling (optimization)}\index{bounds (optimization)}%
Physical parameters are positive and of very different magnitudes: a modulus of
$2\times10^5$\,MPa, an exponent of $5$. Two simple devices remove most of the
difficulties.
\begin{itemize}
  \item \emph{Scaling.} Work with dimensionless variables, $p_j/p_j^{\mathrm{ref}}$ or
  $\log p_j$. The logarithm also enforces positivity and makes a relative change the
  natural unit, which is what the Gauss--Newton eigenvalues of Chapter~\ref{ch:id}
  measure.
  \item \emph{Bounds.} Projected methods keep the iterates in a box: L-BFGS-B
  projects the quasi-Newton step and freezes the variables at their
  bounds~\cite{ByrdLuNocedalZhu1995}; trust-region reflective methods do the same
  for least squares. At a solution on a bound, the gradient is not zero but points
  out of the admissible set: the optimality condition is a variational inequality,
  and the multiplier of the active bound is the corresponding component of the
  gradient, as for a contact force.
\end{itemize}

\begin{table}[h]
\centering\small
\renewcommand{\arraystretch}{1.15}
\begin{tabular}{@{}>{\raggedright}p{3.0cm}>{\raggedright}p{3.4cm}>{\raggedright}p{3.6cm}>{\raggedright\arraybackslash}p{3.6cm}@{}}
\toprule
\sffamily method & \sffamily needs & \sffamily convergence & \sffamily typical use\\
\midrule
steepest descent & gradient & linear, $(\kappa-1)/(\kappa+1)$ & almost never\\
Newton & gradient, Hessian & quadratic, local & small smooth problems\\
BFGS, L-BFGS & gradient (adjoint) & superlinear & many parameters\\
Gauss--Newton & residuals, $\mat{S}$ & quadratic if zero residual & small residuals\\
Levenberg--Marquardt & residuals, $\mat{S}$ (direct diff.) & as Gauss--Newton, safeguarded & few parameters, least squares\\
Nelder--Mead & values & none in general & 2--3 parameters, noisy cost\\
global methods & values & probabilistic & starting points\\
\bottomrule
\end{tabular}
\caption{Minimization algorithms; $\kappa$ is the condition number of the Hessian.}
\label{tab:opt-methods}
\end{table}

%======================================================================
\framepart{Summary}
\begin{gbox*}
\begin{itemize}
  \item \textbf{Descent and line search}: a descent direction and a step satisfying
  the Armijo (and Wolfe) conditions give a stationary point; steepest descent
  zigzags in valleys, with the rate $(\kappa-1)/(\kappa+1)$.
  \item \textbf{Newton and BFGS}: Newton uses the Hessian; BFGS builds the inverse
  Hessian from gradient differences (secant condition), positive definite under
  the Wolfe conditions, superlinear.
  \item \textbf{Least squares}: Gauss--Newton keeps the first-order part
  $\mat{S}\T\mat{S}$ of the Hessian; Levenberg--Marquardt damps it, which is a trust
  region and a Tikhonov regularization of the flat directions.
  \item \textbf{Without derivatives}: Nelder--Mead for two or three parameters;
  model-based and global methods otherwise.
  \item \textbf{Scaling and bounds}: logarithmic variables; projected methods for
  bounds.
\end{itemize}
\end{gbox*}

\framepart{Exercises}
\framesub{Short questions}

\exercise{Steepest descent on a valley}\label{exo:opt-sd}
\index{steepest descent!exercise}%
Let $\mathcal J=\tfrac12(p_1^2+100\,p_2^2)$. (a) Starting from $(100,1)$, compute two
steepest-descent steps with exact line search. (b) Give the factor of decrease of
the error per iteration and the number of iterations for six digits. (c) What does
one Newton step do?
\begin{solution}
(a) The gradient is $(p_1,100p_2)$; along $\vd=-\vect g$ the exact step is
$a=\vect g\cdot\vect g/\vect g\T\mat{A}\vect g$. From $(100,1)$: $\vect g=(100,100)$,
$a=2/101$, $\vp_1=(98.02,-0.980)$; then $\vect g=(98.02,-98.0)$ and
$\vp_2\approx(96.1,0.961)$. The iterates zigzag across the valley and progress
by $2\%$ per step. (b) $\kappa=100$, factor $99/101=0.980$; six digits take about
$\ln10^{-6}/\ln0.98\approx680$ iterations. (c) Newton solves
$\mat{A}\vd=-\vect g$ and lands on the minimum in one step.
\end{solution}

\exercise{The secant condition}\label{exo:opt-secant}
\index{BFGS method!exercise}%
(a) Check that~\eqref{eq:opt-bfgs} satisfies $\mat{H}_{k+1}\vy_k=\vs_k$. (b) Show that
$\mat{H}_{k+1}$ is positive definite if $\mat{H}_k$ is and $\vy_k\cdot\vs_k>0$. (c) For a
quadratic cost with Hessian $\mat{A}$, what is $\vy_k$?
\begin{solution}
(a) $(\mat{I}-\rho\vy\vs\T)\vy=\vy-\rho\vy(\vs\cdot\vy)=\vzero$, so the first term
vanishes, and $\rho\vs(\vs\cdot\vy)=\vs$. (b) For $\vz\neq\vzero$, with
$\vw=(\mat{I}-\rho\vy\vs\T)\vz=\vz-\rho(\vs\cdot\vz)\vy$:
$\vz\T\mat{H}_{k+1}\vz=\vw\T\mat{H}_k\vw+\rho(\vs\cdot\vz)^2\ge0$; both terms vanish only if
$\vw=\vzero$ and $\vs\cdot\vz=0$, which gives $\vz=\vzero$. (c) $\vy_k=\mat{A}\vs_k$:
the secant condition is then exact, and BFGS with exact line searches finds the
minimum of a quadratic in at most $n$ steps.
\end{solution}

\exercise{Limits of the damping}\label{exo:opt-lm}
\index{Levenberg--Marquardt method!exercise}%
In~\eqref{eq:opt-lm} with $\mat{D}=\mat{I}$, show that $\nu\,\vd\to-\mat{S}\T\mat{r}$ as
$\nu\to\infty$ and that $\vd$ tends to the Gauss--Newton step as $\nu\to0$ when
$\mat{S}$ has full rank. Along which eigenvectors of $\mat{S}\T\mat{S}$ is a moderate
damping most effective?
\begin{solution}
Write $\vd=-(\mat{S}\T\mat{S}+\nu\mat{I})^{-1}\mat{S}\T\mat{r}$ in the eigenbasis
$(\lambda_i,\vect v_i)$: the component along $\vect v_i$ is
$-(\vect v_i\cdot\mat{S}\T\mat{r})/(\lambda_i+\nu)$. For $\nu\gg\lambda_i$ it is
$-(\cdot)/\nu$: a short steepest-descent step. For $\nu\to0$ it is $-(\cdot)/\lambda_i$:
Gauss--Newton. A damping $\nu$ between the eigenvalues cuts the components with
$\lambda_i\ll\nu$, the flat directions, and keeps those with $\lambda_i\gg\nu$.
\end{solution}

\framesub{Problems}

\exercise{Six methods on the Rosenbrock valley}\label{exo:opt-rosen}
\index{Rosenbrock function}%
The function of Rosenbrock~\cite{Rosenbrock1960},
$\mathcal J=\tfrac12\norm{\mat{r}}^2$ with $\mat{r}=(10(p_2-p_1^2),\,1-p_1)$, has its minimum
at $(1,1)$ at the bottom of a curved valley. From $(-1.2,1)$, minimize it to
$10^{-6}$ with steepest descent, Newton and Gauss--Newton (all with an Armijo
backtracking), Nelder--Mead, BFGS and Levenberg--Marquardt, and count the
evaluations.
\begin{solution}
Steepest descent needs $14\,446$ iterations and $129\,225$ values of $\mathcal J$: it
zigzags along the curved valley. Newton needs 21 iterations, Gauss--Newton 10
(32 values), BFGS 36 iterations with 47 values and gradients, Levenberg--Marquardt
21 residual and 16 Jacobian evaluations, and Nelder--Mead 117 iterations with 219
values. The methods that use curvature, exact (Newton) or from first derivatives
(Gauss--Newton, BFGS, Levenberg--Marquardt), follow the valley; in two dimensions
the simplex does too, at a moderate price.
\lstinputlisting[firstline=5]{python/examples/appF/exo_minimization.py}
\end{solution}

\begin{chapterbib}{99}
\bibitem{Armijo1966} L.~Armijo, Minimization of functions having Lipschitz continuous first partial derivatives, \emph{Pacific J. Math.} 16 (1966) 1--3.
\bibitem{ByrdLuNocedalZhu1995} R.~H.~Byrd, P.~Lu, J.~Nocedal, C.~Zhu, A limited memory algorithm for bound constrained optimization, \emph{SIAM J. Sci. Comput.} 16 (1995) 1190--1208.
\bibitem{ConnScheinbergVicente2009} A.~R.~Conn, K.~Scheinberg, L.~N.~Vicente, \emph{Introduction to Derivative-Free Optimization}, SIAM, 2009.
\bibitem{DennisSchnabel1983b} J.~E.~Dennis, R.~B.~Schnabel, \emph{Numerical Methods for Unconstrained Optimization and Nonlinear Equations}, Prentice-Hall, 1983 (reprinted by SIAM, 1996).
\bibitem{Fletcher1987} R.~Fletcher, \emph{Practical Methods of Optimization}, 2nd ed., Wiley, 1987.
\bibitem{Kelley1999} C.~T.~Kelley, \emph{Iterative Methods for Optimization}, SIAM, 1999.
\bibitem{LagariasEtAl1998b} J.~C.~Lagarias, J.~A.~Reeds, M.~H.~Wright, P.~E.~Wright, Convergence properties of the Nelder--Mead simplex method in low dimensions, \emph{SIAM J. Optim.} 9 (1998) 112--147.
\bibitem{Levenberg1944} K.~Levenberg, A method for the solution of certain non-linear problems in least squares, \emph{Quart. Appl. Math.} 2 (1944) 164--168.
\bibitem{Marquardt1963} D.~W.~Marquardt, An algorithm for least-squares estimation of nonlinear parameters, \emph{J. SIAM} 11 (1963) 431--441.
\bibitem{More1978} J.~J.~Mor\'e, The Levenberg--Marquardt algorithm: implementation and theory, in G.~A.~Watson (ed.), \emph{Numerical Analysis}, Lecture Notes in Mathematics 630, Springer, 1978, 105--116.
\bibitem{NelderMead1965b} J.~A.~Nelder, R.~Mead, A simplex method for function minimization, \emph{Comput. J.} 7 (1965) 308--313.
\bibitem{NocedalWright2006b} J.~Nocedal, S.~J.~Wright, \emph{Numerical Optimization}, 2nd ed., Springer, 2006.
\bibitem{PressEtAl2007} W.~H.~Press, S.~A.~Teukolsky, W.~T.~Vetterling, B.~P.~Flannery, \emph{Numerical Recipes: The Art of Scientific Computing}, 3rd ed., Cambridge University Press, 2007.
\bibitem{Rosenbrock1960} H.~H.~Rosenbrock, An automatic method for finding the greatest or least value of a function, \emph{Comput. J.} 3 (1960) 175--184.
\bibitem{Wolfe1969} P.~Wolfe, Convergence conditions for ascent methods, \emph{SIAM Rev.} 11 (1969) 226--235.
\end{chapterbib}
